
\begin{hintbox}
\textbf{نکته:}
لیست پلتفرم‌های پشتیبانی شده زیر قابل استنداد نیست و صرفا قصد دارد وسعت لینوکس را نمایش بدهد.
\end{hintbox}

%Note:  The following supported-platforms list is not serious documentation.  
%The point is merely to illustrate the breadth of Linux's reach.  

اگر بطور جدی به پورت‌های لینوکس علاقه مندید می‌توانید دو صفحه
\emph{\url{Xose Vazquez Perez's Linux ports page}}\footnote{\texttt{\url{http://web.archive.org/web/20070813000855/http://www.itp.uni-hannover.de/ports/linux_ports.html}}}
و
\emph{\url{Jerome Pinot's Linux architectures list}}\footnote{\texttt{\url{http://web.archive.org/web/20050308130348/http://ngc891.blogdns.net/kernel/docs/arch.txt}}}
را نیز برسی کنید.  (هردو صفحه مذکور در سال ۲۰۰۵ ناپدید شدند)
چرا که پشتیبانی سخت‌افزاری پیچیده‌تر از عملکرد عمومی \LRE{CPU} است و شامل پشتیبانی از
انواع مختلف \LRE{bus} و مسائل ظریف سخت‌افزاری دیگر می‌شود (به‌ویژه برای پورت‌های لینوکس
\emph{\url{PDA / embedded / microcontroller / router}}\footnote{\texttt{\url{http://www.linuxfordevices.com/}}}
).

%If seriously interested in the subject of Linux ports, please see also 

%\emph{Xose Vazquez Perez's Linux ports page} \tturl{http://web.archive.org/web/20070813000855/http://www.itp.uni-hannover.de/ports/linux_ports.html}
%and
%\emph{Jerome Pinot's Linux architectures list} \tturl{http://web.archive.org/web/20050308130348/http://ngc891.blogdns.net/kernel/docs/arch.txt}

% (static mirrors, as both pages vanished in 2005), if only because 
%hardware support is more complex than just generic CPU functionality, 
%encompassing support for myriad bus variations and other subtle hardware
%issues (especially for 

%\href{http://www.linuxfordevices.com/}{Linux PDA / embedded / microcontroller / router ports}
%).  

\begin{latin}
\begin{itemize}
	\item {\bfseries Diverse \emph{\url{PDA / embedded / microcontroller / router}}\footnote{\texttt{\url{http://www.uclinux.org/ports/}}} devices:}
	\begin{itemize}
		\item Advanced RISC Machines, Ltd. \emph{\url{ARM}}\footnote{\texttt{\url{http://www.arm.linux.org.uk/}}} family (StrongARM SA-1110, XScale, ARM6, ARM7, ARM2, ARM250, ARM3i, ARM610, ARM710,ARM7TDMI, ARM720T, and ARM920T, including Sigma Designs DVD systems using ARM cores)
		\item Analog Devices, Inc.'s \emph{\url{Blackfin DSP}}\footnote{\texttt{\url{http://www.linuxfordevices.com/c/a/News/Another-nativeDSP-Linux-port-this-one-to-ADIs-Blackfin/}}}
		\item Axis Communications \emph{\url{ETRAX series}}\footnote{\texttt{\url{https://en.wikipedia.org/wiki/ETRAX_CRIS}}}
			("CRIS» = Code Reduced Instruction Set RISC architecture)
		\item Elan SC520 and SC300
		\item FreeScale \emph{\url{MC68EN302}}\footnote{\texttt{\url{http://www.freescale.com/files/netcomm/doc/data_sheet/MC68EN302.pdf}}}
		\item Fujitsu \emph{\url{FR-V}}\footnote{\texttt{\url{http://ecos.sourceware.org/hardware.html#FR-V}}}
		\item Hitachi \emph{\url{H8}}\footnote{\texttt{\url{http://www.uclinux.org/pub/uClinux/ports/h8/}}} series
		\item Intel i960
		\item Intel IA32-compatibles (Cyrix MediaGX, STMicroelectronics \emph{\url{STPC}}\footnote{\texttt{\url{http://web.archive.org/web/20070626190436/http://www.stmcu.com/forums-cat-132-6.html}}}, ZF Micro ZFx86)
		\item Matsushita \emph{\url{AM3x}}\footnote{\texttt{\url{http://ecos.sourceware.org/hardware.html#MatsushitaAM3x}}}
		\item MIPS-compatibles (Toshiba TMPRxxxx / TXnnnn, NEC \emph{\url{VR}}\footnote{\texttt{\url{http://www.linux-mips.org/wiki/NEC_VR4100}}} series, Realtek 8181")
		\item Motorola 680x0-based machines (Motorola VMEbus boards, 
			\emph{\url{ISICAD Prisma}}\footnote{\texttt{\url{http://ds.dial.pipex.com/town/way/fr30/}}} machines, and Motorola Dragonball \& 
			\emph{\url{ColdFire}}\footnote{\texttt{\url{http://www.uclinux.org/ports/coldfire/}}} CPUs, and Cisco 2500/3000/4000 series routers)
		\item Motorola embedded \emph{\url{PowerPC}}\footnote{\texttt{\url{http://penguinppc.org/embedded/}}}
			(including MPC / PowerQUICC I, II, III families)
		\item NEC \emph{\url{V850E}}\footnote{\texttt{\url{http://ecos.sourceware.org/tools/linux-v850-elf.html}}}
		\item Renesas Technology (formerly Hitachi) SH3/SH4 (\emph{\url{SuperH}}\footnote{\texttt{\url{http://wiki.debian.org/SH4}}})
		\item Samsung \emph{\url{CalmRISC}}\footnote{\texttt{\url{http://ecos.sourceware.org/hardware.html#CalmRISC}}}
		\item Texas Instruments's 
		\emph{\url{DM64x}}\footnote{\texttt{\url{http://www.linuxfordevices.com/c/a/News/Worlds-first-nativeDSP-Linux-port/}}}
			and \emph{\url{C54x DSP}}\footnote{\texttt{\url{http://www.linuxfordevices.com/c/a/News/Embedded-Linux-distro-supports-TI-DSPbased-digital-media-processors/}}} families
		\item Xilinx \emph{\url{PetaLinux}}\footnote{\texttt{\url{http://www.xilinx.com/tools/petalinux-sdk.htm}}}
			(formerly SoftBlaze, formerly Microblaze) soft processor implemented on Xilinx FPGAs
	\end{itemize}

	\item {\bfseries Intel \emph{\url{8086 / 80286}}\footnote{\texttt{\url{http://elks.sourceforge.net/}}}}.
	\item {\bfseries Intel IA32 family:} i386, i486, Pentium, Pentium Pro,
		Pentium II, Pentium III, Celeron, Xeon, and Pentium IV processors,
		as well as IA32 clones from AMD (386DX/DXL/SL/SLC/SX,
		486DX/DX2/DX4/SL/SLC/SLC2/SLC3/SX/SX2, Elan, K5,
		K6/K6-II/K6-III), Cyrix (386DX/DXL/SL/SLC/SX,
		486DLC/DLC2/DX/DX2/DX4/SL/SLC/SLC2/SLC3/SX/SX2, Cyrix III),
		IDT (Winchip, Winchip 2, Winchip 2A/3),
		IBM (486DX/DX2/DX4/SL/SLC/SLC2/SLC3/SX/SX2),
		NexGen (Nx586), Transmeta (Crusoe),
		TI (486DLC/DLC2), UMC (486SX-S, U5D/U5S),
		VIA (C3 Ezra «CentaurHauls", C3-2 «Nehemiah"),
		and others.
	\item {\bfseries Intel/HP \emph{\url{IA64}}\footnote{\texttt{\url{http://www.ia64-linux.org/}}}:} Trillian, Itanium, Itanium2/McKinley
	\item {\bfseries x86-64 family} including AMD Hammer/Opteron/K8/Athlon64/Turion/Phenom/Phenom II/FX/Fusion and
		Intel Prescott/Nocona/Potomac, Core, Atom, Nehalem, Sandy Bridge and Ivy Bridge
	\item {\bfseries Motorola \emph{\url{68020-68040}}\footnote{\texttt{\url{http://www.linux-m68k.org/}}} series (with MMU)}:
		\emph{\url{m68k Mac}}\footnote{\texttt{\url{http://www.mac.linux-m68k.org/}}},
		Amiga, Atari ST/TT/Medusa/Falcon, HP/Apollo Domain,
		\emph{\url{HP9000/300}}\footnote{\texttt{\url{http://www.tazenda.demon.co.uk/phil/linux-hp/}}}, sun3, and
		\emph{\url{Sinclair Q40}}\footnote{\texttt{\url{http://ftp4.de.freesbie.org/pub/misc/tsx-11/680x0/q40/install/}}}.
	\item {\bfseries Motorola/IBM PowerPC family:} Most
		\emph{\url{PowerMac}}\footnote{\texttt{\url{http://penguinppc.org/mac/}}}
		 (including G3/G4/G5)  / CHRP / PReP / POP, \emph{\url{Amiga PowerUP System}}\footnote{\texttt{\url{http://linux-apus.sourceforge.net/}}}, and IBM PPC64 (AS/400, RS/6000, iSeries,	pSeries, PowerMac G5).
	\item {\bfseries \emph{\url{MIPS}}\footnote{\texttt{\url{http://www.linux-mips.org/}}}:} most SGI, Cobalt Qube,
		\emph{\url{DECStation}}\footnote{\texttt{\url{http://decstation.unix-ag.org/}}}, Sony PlayStation2, and many others
	\item {\bfseries DEC\emph{\url{Alpha}}\footnote{\texttt{\url{http://www.alphalinux.org/}}}}
	\item {\bfseries HP \emph{\url{PA-RISC}}\footnote{\texttt{\url{http://www.parisc-linux.org/}}}}
	\item {\bfseries SPARC International SPARC32 / SPARC64}
	\item {\bfseries Digital \emph{\url{VAX}}\footnote{\texttt{\url{http://vax-linux.org/}}} minicomputers and MicroVAXen}
	\item {\bfseries Mainframes:} 
	\emph{\url{IBM S/390 models G5 and G6 / zSeries models z800, z890, z900, and z990}}\footnote{\texttt{\url{https://www.ibm.com/developerworks/linux/linux390/}}}
		 and Fujitsu AP1000+ (SuperSPARC cluster)
\end{itemize}
\end{latin}

توجه کنید که موارد ذکر شد بعضا تنها یک بار فورک شده اند. در این لیست ریز تا درشت
پورت‌های مختلف لینوکس معرفی شده. در برخی از معماری‌های نادر،
\emph{\url{NetBSD}}\footnote{\texttt{\url{http://www.netbsd.org/}}}
گاها کارآمد تر خواهد بود. همچنان پورت
\emph{\url{Debian GNU/kFreeBSD}}\footnote{\texttt{\url{http://www.debian.org/ports/kfreebsd-gnu/}}}
به‌اندازه کافی سازگار و مناسب میباشد. این پورت مجهز به کد یوزر اسپیس گنو/لینوکس بر روی کرنل فری بی اس دی با پرفرمنس بالا و پایداری بالا،
\emph{\url{OpenIndiana}}\footnote{\texttt{\url{https://en.wikipedia.org/wiki/OpenIndiana}}}
و یا دیگر 
\emph{\url{Illumos distribution}}\footnote{\texttt{\url{https://en.wikipedia.org/wiki/Illumos#Relatives}}}
میباشد و می‌تواند چیز دیگری مشابه کرنل اپن سولاریس فراهم کند.

%Note that some items listed were probably one-time forks, little or not
%at all maintained since creation.  On some of the rarer architectures, \emph{NetBSD} \texttt{http://www.netbsd.org/} may be more practical.(The \emph{Debian GNU/kFreeBSD} \texttt{http://www.debian.org/ports/kfreebsd-gnu/}
% port should also be solid enough to 
%serve as a compromise option, furnishing GNU/Linux userspace code on the
%high performance / high stability FreeBSD kernel, and
%\emph{OpenIndiana} \texttt{https://en.wikipedia.org/wiki/OpenIndiana}
% or another 
%\emph{Illumos distribution} \texttt{https://en.wikipedia.org/wiki/Illumos#Relatives}
% can provide something similar on the OpenSolaris kernel.)

